\diarytitle{0051}{ルジャンドル方程式（$n=2$）の多項式解とルジャンドル多項式 $P_2(x)$}

$(1-x^2)y'' - 2xy' + n(n+1)y = 0$ に $y=\sum a_k x^k$ を代入すると
\[
\sum_{k=0}^{\infty}(k+2)(k+1)a_{k+2}x^{k}
-\sum_{k=0}^{\infty} k(k-1)a_k x^{k}
-2\sum_{k=0}^{\infty} k a_k x^k
+ n(n+1)\sum_{k=0}^{\infty} a_k x^k = 0
\]
（$(1-x^2)y''$ の $x^2 y''$ 部分は添字をずらさずそのまま $k(k-1)a_k x^k$ となることに注意）。$x^k$ の係数をまとめると
\[
(k+2)(k+1)a_{k+2} = \big[k(k-1) + 2k - n(n+1)\big]a_k = \big[k(k+1) - n(n+1)\big]a_k
\]
すなわち
\[
a_{k+2} = \frac{k(k+1) - n(n+1)}{(k+1)(k+2)}\,a_k
\]
が一般の漸化式である。ここに $n=2$（$n(n+1)=6$）を代入すると
\[
a_{k+2} = \frac{k(k+1)-6}{(k+1)(k+2)}\,a_k
\]
$a_1=0$ とすると奇数次の係数はすべて $0$（$a_3=a_5=\cdots=0$）。偶数次は
\[
k=0:\quad a_2 = \frac{0-6}{1\cdot 2}a_0 = -3a_0
\]
\[
k=2:\quad a_4 = \frac{2\cdot 3 - 6}{3\cdot 4}a_2 = \frac{0}{12}a_2 = 0
\]
以降 $a_6=a_8=\cdots=0$ も同様に $0$ となるから、級数は $a_0, a_2$ の2項で止まり
\[
y = a_0 + a_2 x^{2} = a_0\left(1 - 3x^{2}\right)
\]
これが求める多項式解である。$x=1$ で値 $1$ をとるように定数 $a_0$ を定めると
\[
a_0(1-3) = 1 \quad\Longrightarrow\quad a_0 = -\frac{1}{2}
\]
よって
\[
\boxed{\; P_2(x) = -\frac{1}{2}\left(1-3x^{2}\right) = \frac{3x^{2}-1}{2} \;}
\]
（検算: $(1-x^2)y''-2xy'+6y$ に $y=\frac{3x^2-1}{2}$ を代入すると、$y'=3x$, $y''=3$ より
$(1-x^2)\cdot 3 - 2x\cdot 3x + 6\cdot\frac{3x^2-1}{2} = 3-3x^2-6x^2+9x^2-3 = 0$ となり方程式を満たす。また $P_2(1)=\frac{3-1}{2}=1$ で正規化条件も満たす。）
